\cotefile{162-5}
\scanpage{37}
éventuellement muni de la sous-catégorie pleine de tous les systêmes locaux
sur X, et du foncteur cohomologique sur $D^b_{lc}(X)$ à valeurs dans lad\deleted{d}it\deleted{\unsure{u}}e
catégorie, et bien sûr du produit tensoriel (mais alors on sort de $D^b$
pour entrer dans $D^-$ , redactor demerdetur).

II n-catégories, catégories n-uples, et (Gr)-catégories.

Appelons bi-catégori\deleted{\unsure{i}}e une catégorie dans(Cat), i.e. un triple
$(C^{\deleted{o}}_{0},C^{\deleted{1}}_{1},$ \deleted{) où}$p_1,p_2,\pi$ ), où $C^o$ est une catégorie, $p_1,p_2$: $C^{\deleted{1}}_{1}\rightrightarrows C^{\deleted{o}}_{0}$ deux
foncteurs, $\pi:(C^{\deleted{1}}_{1},p_1)x_{C^o}(C^{\deleted{1}}_{1},p_2)\longrightarrow C^{\deleted{1}}_{1}$ un foncteur, avec les axiomes
habituels. Interprétant (Cat) comme la sous-catégorie pleine de (Ssimpl)
des objets qui .... , et de même pour les catégories à valeurs dans une
catégorie telle (Cat), on trouve que la catégorie \inserted{(Bicat)} des (Cat)-caté\deleted{t}gories
est équivalente à la sous-catégorie pleine de la catégorie des ensembles
bi-semisi\deleted{\unsure{p}}mplicia\deleted{y}ux

\[
\begin{array}{llll}
X^{\deleted{oo}}_{00} & X^{\deleted{\unsure{o1|d1}}}_{01} & \ldots & X^{\deleted{o,n}}_{0n}\\
X^{\deleted{1o}}_{10} & X^{\deleted{11}}_{11} & \ldots & X^{\deleted{1,n}}_{1n}\\
\multicolumn{4}{l}{\ldotp\;\ldotp\;\ldotp\;\ldotp\;\ldotp\;\ldotp\;\ldotp}
\end{array}
\]

formée des objets qui ligne par ligne et colonne par colonne satisfoht la
condition d'exctitude familière (transformer sommes amalgamées en produits
fibrés). De même la catégorie \deleted{des} (Tricat) des (Bicat)-catégories
\unsure{\deleted{'}s\deleted{s}'}interprète en t ermes d'objets tri-\inserted{semi}simpliciaux, etc.

Exemple: soit C une catégorie, posons $C^{\deleted{o}}_{0}=C$, $C^{\deleted{1}}_{1}=\underline{\mathrm{Fl}}(C)$, avec structures
$p_1,p_2$, évidentes (d'où une (Cat)-catétorie dont les composantes semi-sim-\\
pliciales supérieures sont les catégories $\underline{\mathrm{Fl}}^n(C)$ des n-flêches de C ...).
L'ensemble bi-semisimplicial correspondant a pour composantes les ensembles
$\mathrm{Hom}(\Delta^n x\Delta^m,C)$. Par exemple, les objets de $X^{\deleted{11}}_{11}$ sont les carrés commutatifs
de flêches de C, qu'on peut interpréter(de deux manières différentes) comme

\margin{36}
